Having shown that alignment reshapes attention intrinsically (\Cref{tab:align}) without moving downstream numbers (\Cref{tab:bench}), we now ask what the three configurations actually do, where in the model they do it, and where the method breaks.

\subsection{A behavioral dissociation between components}
\label{sec:analysis-pheno}
The three adaptation sites do not produce one uniform ``improvement from saliency supervision''; they produce distinct phenotypes. The LM-aligned models (LM-only and LM+projector) earn by far the cleanest saliency (\Cref{tab:align}) yet sit at base-model level downstream (\Cref{tab:bench}), with capabilities preserved. The projector-only model is the opposite: its saliency barely improves over plain caption fine-tuning, but it is the one variant whose downstream behavior shifts sharply---hallucination worsens (POPE $84.3\!\to\!73.2$) and the bias rate falls ($41.7\!\to\!35.9$).

The takeaway is a clean \emph{decoupling}: saliency-map quality does not predict downstream grounding. Counting makes this sharpest. Across components, attention quality and counting accuracy are not merely uncorrelated but \emph{anti}-correlated---training the projector \emph{improves} TallyQA counting while training the language model, which yields by far the best maps, degrades it (\Cref{app:breakdown}). The model that looks best does not answer best. This echoes the attention-overlap-versus-counting null we find across attribution methods (\Cref{app:counting}), and it cautions against treating sharper attention as a proxy for better grounding.

\subsection{Where the adaptation lives}
\label{sec:analysis-where}
Parameter drift between each checkpoint and the base model explains why the LM-aligned variants behave identically and why the projector cannot stand in for the LM. The language model's relative update is essentially \emph{invariant} to whether the projector is also trainable---$7.72\times10^{-3}$ (LM-only) versus $7.67\times10^{-3}$ (LM+projector), agreeing to three significant figures---while under joint training the projector moves only $22\%$ of its solo magnitude ($1.36\times10^{-2}$ vs.\ $6.07\times10^{-2}$). The loss splits its gradient across components, but the language-model solution dominates and the two do not compose: LM+projector is the LM solution, not a superposition. (The projector-only model's large projector drift in turn feeds the frozen LM visual features off its training distribution, the likely source of its POPE/grounding degradation in \Cref{sec:analysis-pheno}.) The vision tower is confirmed frozen, its drift at conversion-noise level across all runs.

\label{sec:analysis-heads}
Within the language model the update is spread across all decoder layers and only \emph{modestly} concentrated---layer-level Gini $\approx 0.10$--$0.13$ for attention projections and $\approx 0.06$ for MLP, where a uniform spread is $0$---so we do not characterize the adaptation as sparse or low-rank. It is somewhat sharper at the head level (Gini up to $0.29$), with a standout head (layer~11 query projection, head~22) that dominates its own layer. The more robust signal is \emph{reproducibility}: between the LM-only and LM+projector runs the top-drifting layers overlap with Jaccard $=1.00$, and the top heads with Jaccard $=1.00$ in $8$ of $10$ analyzed cells, so the optimizer reaches essentially the same solution regardless of the projector. This is consistent with the localization-head phenomenon~\citep{kang2025your, kim2025interpreting}, here engaged by \emph{explicit} spatial supervision. Because the update amplifies a stable, reproducible set of pre-existing directions rather than installing large new ones, we conjecture it is well matched to low-rank LoRA adapters (\Cref{sec:regimes})---a hypothesis we have not yet tested. Two caveats bound this analysis and we return to them in \Cref{sec:limitations}: parameter drift is only a proxy for functional change (we do not causally ablate the candidate heads), and reproducibility is established across configurations, not across random seeds.

\subsection{A failure mode on background ``stuff''}
\label{sec:analysis-stuff}
The one place alignment clearly hurts is broad, amorphous background---``stuff'' classes such as snow and sky (\Cref{fig:qual}, bottom). After fine-tuning, the base model's sensible spreading of saliency across a snowy ground is replaced by a withdrawal of mass onto foreground objects. We attribute this to the \emph{supervision distribution} rather than the objective itself. Grounded captions reference discrete, nameable entities---individual objects and groups---and almost never attach a token to a large background region, so the alignment target $\mathbf{T}_t$ practically never places mass on stuff; because cross-entropy still rewards \emph{mentioning} the background while the alignment term only ever pulls mass toward objects, the model learns to treat ``unannotated'' as ``unattended.'' The control is telling: background entities that \emph{are} annotated because they have definite extent (plants, branches) stay well localized---it is specifically the un-annotatable, region-filling categories that degrade. The fix follows directly: COCONut provides full \emph{panoptic} masks including stuff segments, so extending supervision to stuff tokens---or simply exempting background tokens from the alignment term---would remove the implicit ``attend only to objects'' pressure. We leave stuff-aware supervision to future work.
